\Def Операция f идемпотентна, если $\forall a \rightarrow f(a,\,a)\,=\,a$.

\Def Разреженная таблица или Sparse Table — двумерная структура
данных ST[i][j], построенная на бинарной идемпотентной операции F, для которой выполнено следующее:

\begin{center}
$ST[i][j] = F(A[j],\,A[j + 1],\,\dots,\,A[j + 2^i - 1]),\,i \in [0 \dots log_2\,N]$
\end{center}

$ST[i][j] =
\begin{cases}
F(ST[i-1][j],\,ST[i-1][j+2^{i-1}])\text{, если i > 0}\\
A[j]\text{, если i = 0}
\end{cases}$

Заметим, что в этой таблице хранятся результаты функции F на всех отрезках, длины которых равны степеням двойки.
Выполним сначала предподсчет, суть которого в вычислении массива $floor\_log[i]\,=\,\lfloor log_2 i \rfloor$. Теперь заметим, что для отрека [l, r] верно, что

{\centering $F(A[l],\,A[l + 1],\,\dots,\,A[r]) = F(ST[i][l],\,ST[i][r - 2^i + 1])$,\par}
где $i\,=\,floor\_log [r - l + 1]$

\drawsomelarge{pix/sparse_table.png}

\textbf{Асимптотика}

В таблице хранятся результаты функции F на всех отрезках, длины
которых равны степеням двойки. Однако для каждого $j \in [0 \dots log_2\,N]$
таких отрезков не более N, откуда потребляемая память составит
$O(N\,log\,N)$.\\
Теперь время построения. Заметим, что каждая ячейка
пересчитывается за $O(1)$, откуда время построения $O(N\,log\,N)$.\\
И последнее — ответ на запрос. Заметим, что это всего лишь
вычисление функции от двух значений, что работает за $O(1)$.

\subsubsection*{Область применимости}

Можно решать static RMQ

% RSQ ни в каких постановках задачи нельзя решать, т.к. сумма не идемпотентна.

RSQ за $O(1)$ ни в каких постановках задачи нельзя решать, т.к. сумма не идемпотентна. Однако RSQ можно решить с помощью SparseTable, если разбить запрос суммы на отрезки, длины которых - это степени двоек. Таких отрезков будет $\log_2{n}$ - в силу представления числа в двоичной системе счисления, а ответ для каждого из таких отрезков уже предподсчитан в SparseTable (По определению), значит получаем асимптотику $O(\log n)$ на запрос RSQ.   

\pagebreak